\section{Why the denominators are cubes}
Transfer matrices for strip colourings are classical; see Salas and
Sokal \cite[Section 3]{salassokal}. Here the defect budget gives an
additional denominator bound.
Let a slice have $s$ sites and $q\geq1$ colours. Let $h(r)$ and
$b(p,r)$ be fixed nonnegative integer costs within a slice and between
consecutive slices. For $d\geq0$, let $L_d(\ell)$ count labelled
strings of $\ell\geq1$ slices with total cost exactly $d$.
Put $T_0(p,r)=\ind_{h(r)+b(p,r)=0}$ and $D(z)=\det(I-zT_0)$.
\begin{theorem}\label{thm:budget}
The reduced denominator of $\sum_{\ell\geq1}L_d(\ell)z^\ell$ divides
$D(z)^{d+1}$. If a group acting on slices preserves $h$ and $b$, the
same assertion holds with the orbit transition matrix replacing $T_0$.
\end{theorem}
\begin{proof}
Use states $(r,e)$ with remaining budget $0\leq e\leq d$, initial
weight $\ind_{e=d-h(r)}$, and terminal weight $\ind_{e=0}$.
The transition decreases $e$ by $h(r)+b(p,r)$. Ordering states by
budget makes the matrix block triangular, with $d+1$ identical
diagonal blocks $T_0$. Thus $\det(I-zM)=D(z)^{d+1}$.
The path interpretation gives the generating function
$z w_0(I-zM)^{-1}v$; the adjugate formula proves the assertion.
For an orbit $C$, the number of $r\in C$ with each prescribed cost
is independent of the representative $p$, because the group preserves
costs. Use these numbers as transition weights and sum initial weights
over orbits. Induction preserves total weight in every orbit and
budget. The diagonal blocks are again identical, proving the assertion.
\end{proof}

For the present arrays, use the group generated by colour permutations
and slice reversal. Represent each orbit by the lexicographically
smaller of the two row-major canonical labellings of the slice and its
reversal. Reversal groups transfer states; it does not identify arrays
in the final count. For fixed width $k$, use the costs in Section 2.
For fixed height $n$, use column slices and
\[
h(r)=0,\qquad b(p,r)=\sum_{i=1}^{n}\ind_{p_i=r_i}
+\sum_{i=1}^{n-1}(\ind_{p_i=r_{i+1}}+\ind_{p_{i+1}=r_i}).
\]
The first sum counts horizontal bonds. Denote the orbit zero-cost
matrix by $B$, and put $D=\det(I-zB)$. For fixed width, columns
indexed by slices with $h(r)>0$ are zero in $B$; deleting their rows
and columns leaves this determinant unchanged. These rules specify
every entry without a preassigned transition table.

\section{All displayed companion recurrences}
Let $T(n,k)=a_k(n)$. Use $x=n$ in a fixed column and $x=k$ in a
fixed row, and set each sequence to zero at index zero.
For $Q(z)=D(z)^3=\sum_{j=0}^{r}q_jz^j$, with $q_0=1$, its
associated recurrence means
\begin{equation}\label{eq:allrec}
f(x)=-\sum_{j=1}^{r}q_j f(x-j).
\end{equation}
\begin{corollary}\label{cor:all}
Every displayed empirical recurrence in A229637, Formula, revision 8
(27 April 2021), holds. In the following table $r$ is the least
eventual rational constant-coefficient order, and $N$ is the sharp
starting index of \eqref{eq:allrec}. The finite vector identity below
has index $N_v$. The last column is the residual at $x=N-1$.
\end{corollary}
\begin{center}\small
\begin{tabular}{lrrrr}
\toprule
Sequence&$r$&$N$&$N_v$&Residual\\
\midrule
$T(x,2)$&3&6&6&$-2$\\
$T(x,3)$&6&8&8&$-88$\\
$T(x,4)$&6&10&11&25\\
$T(x,5)$&18&21&21&$-80$\\
$T(x,6)$&27&31&32&128\\
$T(x,7)$&57&61&62&$-20736$\\
\midrule
$T(2,x)$&6&7&7&$-20$\\
$T(3,x)$&9&13&13&8\\
$T(4,x)$&18&22&22&$-1280$\\
$T(5,x)$&30&34&34&8211456\\
$T(6,x)$&69&73&73&88260480000\\
\bottomrule
\end{tabular}
\end{center}
\begin{proof}
Use the orbit transfer above with budget two, initializing
$w_0(C,e)=|C|\ind_{e=2-h(C)}$. Its terminal weight is the labelled
array count $L$. A constant $n\times k$ array has $(k-1)(3n-2)$
mistakes. Thus the canonical count is $L/6$, except at $(1,3)$,
where it is $(L+3)/6=1$. All other counted arrays have colour orbits
of size six, as in Section 2.

For each table row, the exact finite vector identity is
\[
\sum_{j=0}^{r}q_j w_{N_v-1-j}=0.
\]
The accompanying \texttt{verify\_siblings.py} constructs all orbit
transitions from the three-colour slices, computes $D=\det(I-zB)$
by integer characteristic-polynomial arithmetic, and verifies this
identity in every coordinate. Subsequent transfer powers give
\eqref{eq:allrec} for $x\geq N_v$.
The exceptional constant array $(1,3)$ occurs in none of these identities.
Where $N_v=N+1$, the scalar residual at $N$ is also zero, establishing
the stated range. These are finite vector identities rather than an
extrapolation from an initial list of terms.

The file \texttt{sibling-output.json} supplies $D,Q,P$ in ascending
coefficient order, where
$P_x=\sum_{j=0}^{\min(x,r)}q_jf(x-j)$ for $x<N$.
It verifies the displayed $P_{N-1}$; subsequent coefficients vanish
by the recurrence. For every row, the Euclidean algorithm modulo
$101$ ends in a nonzero constant, with neither leading coefficient
zero. Thus $\gcd(P,Q)=1$ over $\mathbb Q$.
As in Corollary~\ref{cor:minimal}, the denominator has minimal degree
$r$, and $P_{N-1}\ne0$ makes $N$ sharp.
\end{proof}

For orientation, the first four fixed-column polynomials are
\[
\begin{array}{ll}
k=2:&D=1-3z,\\
k=3:&D=1-5z+2z^2,\\
k=4:&D=1-7z+4z^2,\\
k=5:&D=1-14z+45z^2-15z^3-36z^4+19z^5-2z^6.
\end{array}
\]
For row two, $D=1-2z-2z^2$; cubing it gives exactly the six
coefficients in A229637. For columns four, six and seven and rows
three through six, that entry specifies only orders, not coefficients;
the determinant construction specifies the recurrences explicitly.

\begin{corollary}\label{cor:barker}
The two-column sequence is A229600, and
\[
\sum_{n\geq1}a_2(n)z^n
=\frac{z^2(6-15z+13z^2-2z^3)}{(1-3z)^3},\qquad
a_2(n)=3^{n-5}(32n^2+31n-30)\quad(n\geq3).
\]
These are Barker's conjectures in A229600, Formula, 16 March 2018.
They also follow from Shrock's earlier open-ladder partition function
\cite[Eqs.~(5.1)--(5.6)]{shrock} by setting $q=3$, $v=t-1$ and
taking the coefficient of $t^2$: precisely,
$\sum_{n\geq1}a_2(n)z^n=z^2[t^2]\Gamma(S,3,t-1,z)/6$.
The empirical generating functions credited to Mathar in A229632
(29 September 2013) and Barker in A229633 (20 September 2018)
also hold \cite{widthtwo,widththree,widthfour}.
\end{corollary}
\begin{proof}
Swap the two cells in each even-numbered row. Horizontal and crossed
diagonal bonds become the horizontal and vertical bonds of A229600's
ladder. This colour-equivariant graph isomorphism gives a bijection
on canonical colour orbits, even though individual row-major labels
may change. The numerator in Corollary~\ref{cor:all} is
$6z^2-15z^3+13z^4-2z^5$.
Using $[z^m](1-3z)^{-3}=\binom{m+2}{2}3^m$ yields the stated
formula for $n\geq5$; it also gives $39,202$ at $n=3,4$.
For width three, its numerator is
$zD^3-4z^2(2z-5)(12z^4-45z^3+47z^2-17z+2)$.
For width four, it is
$z(3+59z-976z^2+4803z^3-9539z^4+7326z^5-1825z^6-76z^7+25z^8)$.
These are precisely the two cited generating functions.
\end{proof}

The two-row case has a separate exact antecedent. Specialize
Yurishchev's block \cite[Eqs.~(5), (8), (9)]{yurishchev} to three
colours, give the four intercolumn bonds weight $t$, and set all
other couplings to zero. The normalized equal-colour and distinct-colour
slice sums give
\[
H(t)=\begin{pmatrix}
2+t^4&\sqrt2(1+2t^2)\\
\sqrt2(1+2t^2)&4t+2t^2
\end{pmatrix},\qquad b=\begin{pmatrix}\sqrt3\\\sqrt6\end{pmatrix}.
\]
For example, from an equal-colour slice there is one equal-colour
successor of weight $t^4$ and two of weight one. For a fixed
equal-colour slice, two distinct-colour successors have weight one and
four have weight $t^2$; normalization divides their sum by $\sqrt2$.
From a distinct-colour slice,
four distinct successors have one equal bond and two have two.
The free-boundary labelled polynomial is $b^{\mathsf T}H(t)^{k-1}b$:
the all-ones boundary vector has coordinates $b$ in this normalized
invariant basis. Summing the geometric series and taking its $t^2$
coefficient gives
\[
\sum_{k\geq1}T(2,k)z^k
=\frac16[t^2]z b^{\mathsf T}(I-zH(t))^{-1}b
=\frac{2z^2(1-z)(3+5z-36z^2+10z^3)}{(1-2z-2z^2)^3}.
\]
This is also Barker's empirical generating function in A229638,
Formula, 20 September 2018 \cite{rowtwo}. The free-boundary
contraction and coefficient extraction above make the connection
explicit; the published block, rather than this particular finite
strip formula, is Yurishchev's result.

\begin{corollary}\label{cor:rowone}
One has $T(n,1)=0$ for $n\geq1$ and
\[
T(1,k)=\binom{k-1}{2}2^{k-4}\quad(k\geq4),\qquad
(T(1,1),T(1,2),T(1,3))=(0,0,1).
\]
Consequently $T(1,k)=6T(1,k-1)-12T(1,k-2)+8T(1,k-3)$ for
$k\geq7$, with residual $-4$ at $k=6$ and minimal eventual order three.
For $k\geq4$ this is Sloane's A052482$(k-2)$ \cite{sloane}.
In this notation the reciprocal sums recorded there by Amiram Eldar
(28 September 2025) are
\[
\sum_{k\geq4}\frac1{T(1,k)}=6-8\log2,\qquad
\sum_{k\geq4}\frac{(-1)^k}{T(1,k)}=10-24\log(3/2).
\]
\end{corollary}
\begin{proof}
Width one has no bonds. For one row of width $k\geq4$, choose two
equal bonds among $k-1$, then choose the first colour in three ways
and each unequal bond in two ways. Divide by six. The smaller values
follow directly. The third difference of $T(1,k)/2^k$ is zero for
$k\geq7$; its nonzero quadratic term makes the reduced denominator
$(1-2z)^3$, so the order is minimal. Substitution gives the residual.
For the sums, put $m=k-2$ and subtract the logarithmic series for
$1/m$ and $1/(m+1)$ to obtain
\[
S(x)=\sum_{m\geq2}\frac{x^m}{m(m+1)}
=1-\frac{x}{2}+\left(\frac1x-1\right)\log(1-x),\quad 0<|x|<1.
\]
Evaluate $8S(1/2)$ and $8S(-1/2)$.
\end{proof}

\begin{corollary}\label{cor:zero}
Let $U(n)$ be the number of labelled $n\times5$ arrays in A232917:
no horizontal or diagonal neighbours have values summing to two.
Then
\[
\sum_{n\geq1}U(n)z^n
=\frac{6z(4-7z+z^2)(2-8z-5z^2+3z^3)}
 {1-14z+45z^2-15z^3-36z^4+19z^5-2z^6}.
\]
This proves Barker's empirical generating function (6 October 2018)
and Hardin's order-six recurrence in A232917 \cite{zerowidthfive},
the latter for $n\geq7$.
\end{corollary}
\begin{proof}
Replace each colour $x$ by $2-x$ in alternate columns. Every bond
changes column parity, so a complementary pair becomes an equal pair,
and conversely. This bijects the arrays with the labelled zero-defect
arrays of Section 2. The zero-cost orbit matrix $B$ has six states,
so $z u(I-zB)^{-1}\mathbf1$ has denominator $D$ and numerator of
degree at most six, where $u(C)=|C|$ for proper slices.
The first six zero-defect counts, calculated by the same transitions,
give the displayed numerator. The accompanying checker verifies its
coefficients exactly. Expanding $D$ gives the stated recurrence.
\end{proof}

\begin{center}\small
\begin{tabular}{>{\raggedright\arraybackslash}p{.43\linewidth}>{\raggedright\arraybackslash}p{.49\linewidth}}
\toprule
Source item&Coverage\\
\midrule
A229634, Formula; A229637, column 5&Proved; sharp range and minimal order.\\
A229637, columns 2 and row 2; A229600 and A229638&Reproved from prior
formulas: Shrock (2000), Yurishchev (2001); explicit extractions above.\\
A229637, remaining displayed columns and rows&Proved;
Corollaries~\ref{cor:all} and \ref{cor:rowone}.\\
A229632 and A229633, Formula&Hardin's recurrences and Mathar's and Barker's
empirical generating functions proved.\\
A052482, Name; Barker's generating function and recurrence&Prior formulas cited
(Sloane, 2000; Barker, 2017); one-row interpretation proved.\\
A052482, reciprocal sums&Prior formulas cited (Eldar, 2025), with the deductions above.\\
A232917, empirical recurrence and Barker generating function&Proved;
Corollary~\ref{cor:zero}; its denominator is the five-column zero-defect determinant.\\
A052482, folded-cube interpretations; A229606, other widths&Not covered.\\
\bottomrule
\end{tabular}
\end{center}
